\section{Conclusion}
\label{sec:conc}

This work examined whether SIMD optimization principles developed for AIMer v2
remain useful after the transition from AIM2 to AIM3.  We implemented the same
targeted scope with AVX2 and AVX-512: SHAKE x1/x4, party-batched field
arithmetic, binary matrix accumulation, and operation-wise MPC scheduling.
The implementation retains every AIM3-specific equation, constant, random-byte
order, and serialization rule.  Its 1,800 byte-exact KAT responses and direct
differential tests agree with the reference backend.

On the evaluated Tiger Lake processor, AVX2 accelerates signing and
verification by $3.21$--$4.41\times$ over reference, depending on the operation
and parameter set.  AVX-512 raises the improvement to $4.59$--$6.55\times$ and
is $1.43$--$1.52\times$ faster than the matched AVX2 implementation for these
operations.  The result supports our central claim: the low-level SIMD
principles are portable across AIMer generations even though the
algorithm-specific AIM2 code is not.

The study is limited to one AVX-512 microarchitecture and a targeted rather
than full-scheme SIMD scope.  Phase~2/3 field operations and tree traversal are
not reorganized into wider batches, and constant-time behavior has not been
formally verified.  The adaptive kernel measurements also require one final
interrupt-free repetition before they can be treated as publication-ready.
Future work will evaluate additional Intel and AMD microarchitectures, complete
that kernel validation, and investigate the remaining scalar regions without
changing the AIM3 transcript.
